% GENERATED FILE. Edit canonical lessons, then run npm run generate:books.
% canonical-source-hash: fnv1a64:3e513d5f2655f718
% canonical-lessons: UR-C281-thagna, UR-C281-dhoka, UR-C281-rishvat, UR-C281-badmash, UR-C281-jali

\chapter{Cheats}
\label{ch:ur-cheats}
\hlchaptermodality{281}

\begin{chapteropening}
\textbf{By the end of this chapter:} \emph{I can say how someone was cheated and what to do if it happens.}

Complete the last lesson of chapter 281: I can say how someone was cheated and what to do if it happens.
\end{chapteropening}

\section[\texorpdfstring{\ur{ٹھگنا} (ṭhagnā)}{ٹھگنا (ṭhagnā)}]{\ur{ٹھگنا} (ṭhagnā) — to cheat; to swindle}

\label{lesson:UR-C281-thagna}



\begin{quote}
Before the new one: say the Urdu for to tie up, then the Urdu for a looter.
\end{quote}

\subsection*{You'll want to know: \ur{ٹھگنا}}
\textbf{\ur{ٹھگنا}} — \emph{ṭhagnā} — \textquotedblleft{}to cheat; to swindle\textquotedblright{}.

\begin{quote}
\ur{اس} \ur{نے} \ur{مجھے} \ur{ٹھگا۔}
\end{quote}

\begin{quote}
\emph{us ne mujhe ṭhagā.}
\end{quote}

\begin{quote}
He cheated me.
\end{quote}

Past, with \ur{نے} (he or she …): \ur{اس} \ur{نے} \ur{دکاندار} \ur{کو} \ur{ٹھگا}, \emph{us ne dukāndār ko ṭhagā}.

\begin{cousinweb}
\textbf{\ur{ٹھگنا}} (\emph{ṭhagnā}) is inherited from \emph{*ṭhagg}, a form CDIAL reconstructs (CDIAL 5489), which also lists the noun \emph{ṭhag}, a swindler. English \emph{thug} was taken from that noun.
\end{cousinweb}

\begin{grammarlens}[title={What you've built}]
One more word for telling what happened.
\end{grammarlens}

\subsection*{Your turn}
\begin{itemize}
\raggedright
  \item \emph{Say it:} \emph{ṭhagnā}
  \item \emph{Say it:} \emph{ṭhagnā}, once more
  \item \emph{Say it:} \emph{us ne mujhe ṭhagā.} — He cheated me.
  \item \emph{Recall:} say the Urdu for to tie up, then the Urdu for a looter, then say \emph{ṭhagnā} again
\end{itemize}

\begin{tcolorbox}[breakable,colback=teal!4,colframe=teal!35!black,title={Before you move on}]
What does \ur{ٹھگنا} mean? (\textquotedblleft{}To cheat; to swindle\textquotedblright{}.) Say it once more.
\end{tcolorbox}
\section[\texorpdfstring{\ur{دھوکا} (dhokā)}{دھوکا (dhokā)}]{\ur{دھوکا} (dhokā) — a deception; a trick}

\label{lesson:UR-C281-dhoka}



\begin{quote}
Before the new one: say the Urdu for a looter, then the Urdu for to cheat.
\end{quote}

\subsection*{You'll want to know: \ur{دھوکا}}
\textbf{\ur{دھوکا}} — \emph{dhokā} — \textquotedblleft{}a deception; a trick\textquotedblright{}.

\begin{quote}
\ur{اگر} \ur{یہ} \ur{دھوکا} \ur{ہے} \ur{تو} \ur{پولیس} \ur{بلائیے۔}
\end{quote}

\begin{quote}
\emph{agar yih dhokā hai to polīs bulāiye.}
\end{quote}

\begin{quote}
If this is a trick, call the police.
\end{quote}

Masculine.

\begin{cousinweb}
\textbf{\ur{دھوکا}} (\emph{dhokā}) is inherited: CDIAL lists the Hindi cousin \emph{dhokā}, \textquotedblleft{}deceit\textquotedblright{}, under \emph{*dhrōkṣa}, a form it reconstructs (CDIAL 6894). \textbf{\ur{دھوکا} \ur{دینا}} is to deceive someone.
\end{cousinweb}

\begin{grammarlens}[title={What you've built}]
One more word for telling what happened.
\end{grammarlens}

\subsection*{Your turn}
\begin{itemize}
\raggedright
  \item \emph{Say it:} \emph{dhokā}
  \item \emph{Say it:} \emph{dhokā}, once more
  \item \emph{Say it:} \emph{agar yih dhokā hai to polīs bulāiye.} — If this is a trick, call the police.
  \item \emph{Recall:} say the Urdu for a looter, then the Urdu for to cheat, then say \emph{dhokā} again
\end{itemize}

\begin{tcolorbox}[breakable,colback=teal!4,colframe=teal!35!black,title={Before you move on}]
What does \ur{دھوکا} mean? (\textquotedblleft{}A deception; a trick\textquotedblright{}.) Say it once more.
\end{tcolorbox}
\section[\texorpdfstring{\ur{رشوت} (rishvat)}{رشوت (rishvat)}]{\ur{رشوت} (rishvat) — a bribe}

\label{lesson:UR-C281-rishvat}



\begin{quote}
Before the new one: say the Urdu for to cheat, then the Urdu for a deception.
\end{quote}

\subsection*{You'll want to know: \ur{رشوت}}
\textbf{\ur{رشوت}} — \emph{rishvat} — \textquotedblleft{}a bribe\textquotedblright{}.

\begin{quote}
\ur{داروغے} \ur{نے} \ur{رشوت} \ur{مانگی۔}
\end{quote}

\begin{quote}
\emph{dāroghe ne rishvat māṅgī.}
\end{quote}

\begin{quote}
The inspector asked for a bribe.
\end{quote}

Feminine.

\begin{cousinweb}
\textbf{\ur{رشوت}} (\emph{rishvat}) is usually traced, through Persian, to Arabic \emph{rishwa}, \textquotedblleft{}a bribe\textquotedblright{}. It is feminine, so the verb agrees: \textbf{\ur{رشوت} \ur{مانگی}}.
\end{cousinweb}

\begin{grammarlens}[title={What you've built}]
One more word for telling what happened.
\end{grammarlens}

\subsection*{Your turn}
\begin{itemize}
\raggedright
  \item \emph{Say it:} \emph{rishvat}
  \item \emph{Say it:} \emph{rishvat}, once more
  \item \emph{Say it:} \emph{dāroghe ne rishvat māṅgī.} — The inspector asked for a bribe.
  \item \emph{Recall:} say the Urdu for to cheat, then the Urdu for a deception, then say \emph{rishvat} again
\end{itemize}

\begin{tcolorbox}[breakable,colback=teal!4,colframe=teal!35!black,title={Before you move on}]
What does \ur{رشوت} mean? (\textquotedblleft{}A bribe\textquotedblright{}.) Say it once more.
\end{tcolorbox}
\section[\texorpdfstring{\ur{بدمعاش} (badmāsh)}{بدمعاش (badmāsh)}]{\ur{بدمعاش} (badmāsh) — a crook; a rascal}

\label{lesson:UR-C281-badmash}



\begin{quote}
Before the new one: say the Urdu for a deception, then the Urdu for a bribe.
\end{quote}

\subsection*{You'll want to know: \ur{بدمعاش}}
\textbf{\ur{بدمعاش}} — \emph{badmāsh} — \textquotedblleft{}a crook; a rascal\textquotedblright{}.

\begin{quote}
\ur{وہ} \ur{آدمی} \ur{بدمعاش} \ur{ہے۔}
\end{quote}

\begin{quote}
\emph{voh ādmī badmāsh hai.}
\end{quote}

\begin{quote}
That man is a crook.
\end{quote}

Masculine or feminine, by the person meant.

\begin{cousinweb}
\textbf{\ur{بدمعاش}} (\emph{badmāsh}) is usually traced to Persian \emph{bad}, \textquotedblleft{}bad\textquotedblright{}, and Arabic \emph{maʿāsh}, \textquotedblleft{}a living\textquotedblright{}: one who lives by bad means.
\end{cousinweb}

\begin{grammarlens}[title={What you've built}]
One more word for telling what happened.
\end{grammarlens}

\subsection*{Your turn}
\begin{itemize}
\raggedright
  \item \emph{Say it:} \emph{badmāsh}
  \item \emph{Say it:} \emph{badmāsh}, once more
  \item \emph{Say it:} \emph{voh ādmī badmāsh hai.} — That man is a crook.
  \item \emph{Recall:} say the Urdu for a deception, then the Urdu for a bribe, then say \emph{badmāsh} again
\end{itemize}

\begin{tcolorbox}[breakable,colback=teal!4,colframe=teal!35!black,title={Before you move on}]
What does \ur{بدمعاش} mean? (\textquotedblleft{}A crook; a rascal\textquotedblright{}.) Say it once more.
\end{tcolorbox}
\section[\texorpdfstring{\ur{جعلی} (jālī)}{جعلی (jālī)}]{\ur{جعلی} (jālī) — forged; counterfeit}

\label{lesson:UR-C281-jali}



\begin{quote}
Before the new one: say the Urdu for a bribe, then the Urdu for a crook.
\end{quote}

\subsection*{You'll want to know: \ur{جعلی}}
\textbf{\ur{جعلی}} — \emph{jālī} — \textquotedblleft{}forged; counterfeit\textquotedblright{}.

\begin{quote}
\ur{یہ} \ur{نوٹ} \ur{جعلی} \ur{ہے۔}
\end{quote}

\begin{quote}
\emph{yih noṭ jālī hai.}
\end{quote}

\begin{quote}
This banknote is forged.
\end{quote}

\begin{cousinweb}
is usually traced, through Persian, to Arabic \emph{jaʿl}, \textquotedblleft{}a making, a fabrication\textquotedblright{}, with the adjective ending \emph{-ī}. \textbf{\ur{نقلی}} is any copy; a \textbf{\ur{جعلی}} note is one made to pass as real.
\end{cousinweb}

\begin{grammarlens}[title={What you've built}]
One more word for telling what happened.
\end{grammarlens}

\subsection*{Your turn}
\begin{itemize}
\raggedright
  \item \emph{Say it:} \emph{jālī}
  \item \emph{Say it:} \emph{jālī}, once more
  \item \emph{Say it:} \emph{yih noṭ jālī hai.} — This banknote is forged.
  \item \emph{Recall:} say the Urdu for a bribe, then the Urdu for a crook, then say \emph{jālī} again
\end{itemize}

\begin{tcolorbox}[breakable,colback=teal!4,colframe=teal!35!black,title={Before you move on}]
What does \ur{جعلی} mean? (\textquotedblleft{}Forged; counterfeit\textquotedblright{}.) Say it once more.
\end{tcolorbox}
